\section{Проверка на прототипе и ограничения}
\label{sec:limits}

\textbf{Прототип.} Мы проверили обмен памятью модели на прототипе, построенном на llama.cpp \cite{llamacpp}, с
маленькой моделью Qwen3.5-0.8B той же архитектуры, что Qwen3.8-27B, на процессоре.
\begin{itemize}
\item Клиент, который вставил дельту и сам посчитал последний токен порции, генерирует те же 64 токена, что и
клиент, прочитавший порцию локально с той же разбивкой на пакеты, на глубинах 512, 2048 и 8192.
\item Против чтения всего промпта одним пакетом распределения следующего токена расходятся не больше чем на 0.00204
ната на токен (среднее расхождение Кульбака--Лейблера при подстановке верных токенов), а самый вероятный токен
совпадает в 63 или 64 позициях из 64. Так ведёт себя и чисто локальный запуск с другой разбивкой на пакеты, и к
переносу памяти это не относится.
\item Вставка дельты занимает 4--5 миллисекунд и не растёт с глубиной от 512 до 8192 токенов.
\item Пятнадцать видов испорченных дельт клиент отклоняет, и его кэш остаётся побайтно прежним.
\item Журнал сессии сокращает число повторно прочитанных токенов в разговоре из трёх ходов с 1530 до 524.
\end{itemize}
Связку сервера на CUDA с клиентом на Metal и модель 27B на глубине сотен тысяч токенов прототип не проверял. Код
прототипа открыт: \url{https://github.com/relux-works/hybrid-inference}.

\textbf{Доверие к поставщику.} Присланная память модели управляет генерацией: подмена сегмента KV-кэша меняет ход
генерации, не трогая видимый промпт \cite{ganesh2025historyswap,zhu2026r2ecap,zhang2026hijackkv}. В облачном инференсе
клиент проверяет сервер только с его помощью \cite{ong2025toploc}. В гибриде у клиента та же модель, и он может в
простое выборочно пересчитать присланную дельту сам, начиная с рекуррентного состояния, которое заменяется целиком
при каждом оффлоаде. При одних и тех же весах, детерминированных ядрах и одной разбивке на пакеты сверка точная: на
прототипе числа совпали бит в бит. Если сервер считает на CUDA, а клиент на Metal, сверка возможна только с допуском,
и подмену в пределах допуска она не отличит от погрешности. Выборочный пересчёт находит подмену после ответа и
гарантий не даёт. Конфиденциальности схема не обеспечивает: поставщик prefill видит все токены сессии.

\textbf{Ограничения.} Замеры на клиенте сделаны на одной машине, M3 Max; скорости остальных клиентов и сервера оценены
по пиковой производительности и полосе памяти, а доля пика, которую серверный движок использует на prefill, задана
допущением. Движок воспроизводит средние значения и распределения вызовов профиля, но очередей, одновременного
пробуждения многих сессий и пиковой нагрузки не моделирует, и задержки в расчётах ожидаемые. Хранение копий сессий
задано оценкой цены памяти и политикой хранения, которую протокол допускает, а прототип не реализует. Профиль снят с
одной группы пользователей, в основном разработки. Качество генерации мы не проверяли, поэтому сравниваем цену
вызова: если открытая модель чаще ошибается или пишет длиннее, цена выполненной задачи будет другой. Если клиентские
движки ускорят чтение контекста в разы, выигрыш схемы сократится до границы из раздела~\ref{sec:channel}. В обратную
сторону действует работа агентов без оператора круглые сутки: она поднимает загрузку клиента, и в этом отношении
расчёты занижают выгоду гибрида. Следующая версия проверит многоходовую сессию Qwen3.8-27B между сервером на CUDA и
клиентом на Metal с контекстом до 256 тысяч токенов, качество генерации на агентских задачах и полную задержку вызова
с очередью при нескольких клиентах одновременно.
